\chapter{One-Loop Effective Potential in Non-Abelian Gauge Theories}
\label{ch:non-abelian-one-loop}

The mechanism of spontaneous symmetry breaking via radiative corrections,
which we saw explicitly in massless scalar electrodynamics, works equally
well for non-Abelian gauge theories. The computation of the one-loop
effective potential is conceptually the same as in the scalar case,
though the technical details are more involved due to the richer field
content. We follow the treatment of Brandenberger closely.

\section{The Theory}

Consider a non-Abelian gauge theory with gauge group $G$ and generators
$\tau_a$, a set of Dirac bispinor fields $\psi^a$, and a set of scalar
fields $\varphi_i$, each forming a representation space for $G$. The
Lagrangian is
\begin{equation}
    \mathcal{L}(A,\psi,\varphi) = -\frac{1}{4}\mathrm{Tr}(F_{\mu\nu}F^{\mu\nu})
    + i\,\mathrm{Tr}(\bar\psi\slashed{D}\psi)
    + \mathrm{Tr}(\bar\psi M\psi)
    - \mathrm{Tr}(\bar\psi\Gamma\varphi\psi)
    + \frac{1}{2}\mathrm{Tr}(D_\mu\varphi)^\dagger D^\mu\varphi
    - U(\varphi),
    \label{eq:non-abelian-lagrangian}
\end{equation}
where the gauge field, field strength, and covariant derivative are
\begin{align}
    A_\mu &= A_\mu^a\tau_a, \\
    F_{\mu\nu} &= \partial_\mu A_\nu - \partial_\nu A_\mu
    + ig[A_\mu, A_\nu], \\
    D_\mu &= \partial_\mu + igA_\mu.
\end{align}
In $D_\mu$, the $\tau_a$ are the matrices representing the generators
in the appropriate representation space. The trace refers to the
appropriate representation space: if the Higgs and Fermi fields are
represented as column vectors, the trace reduces to the scalar product,
e.g.\ $\mathrm{Tr}(D_\mu\varphi)^\dagger D^\mu\varphi
= (D_\mu\varphi)^\dagger D^\mu\varphi$. The matrix $\Gamma$ denotes the
Yukawa coupling constants.

The Feynman rules in general involve extra fields, the ghost fields
$\omega$ and $\bar\omega$, scalar fields obeying Fermi statistics
(Faddeev and Popov, 1967). They arise as a consequence of the
field-dependent functional determinant in the Faddeev-Popov ansatz for
the generating functional. In covariant gauges, the ghost term in the
effective Lagrangian is $\mathrm{Tr}\,\bar\omega\,\partial_\mu A^\mu\omega$,
and contains no Higgs couplings to ghost fields. Hence to one-loop order
no ghosts arise in our calculation of the effective potential. The
one-loop effective potential therefore takes the form
\begin{equation}
    V_{\text{eff}}(\bar\varphi) = U(\bar\varphi)
    + V_{\text{c.t.}}^{(1)}
    + V_{1\text{-loop}},
    \label{eq:Veff-decomposition}
\end{equation}
where $V_{\text{c.t.}}^{(1)}$ contains all counterterms to order $\hbar$
and $V_{1\text{-loop}}$ is the sum of all one-loop graphs with external
$\varphi$ lines at zero momenta. By charge conservation, the particle
type cannot change within a single loop, so the contributing one-loop
diagrams are those with Higgs bosons, fermions, or gauge bosons running
in the loop. The corresponding contributions are denoted $V_S^{(1)}$,
$V_f^{(1)}$, and $V_g^{(1)}$.

\section{Diagonalization of the Scalar Mass Matrix}

The scalar sector of the theory is described by the Hessian of the
potential, evaluated at the background field $\bar\varphi$,
\begin{equation}
    W^{ij}(\bar\varphi) = \left.\frac{\partial^2 U(\varphi)}
    {\partial\varphi_i\,\partial\varphi_j}\right|_{\varphi=\bar\varphi}.
    \label{eq:Wij-def}
\end{equation}
This matrix controls the quadratic fluctuations around the background.
To see this explicitly, expand around the background
$\varphi_i(x) = \bar\varphi_i + \eta_i(x)$, where $\eta_i$ are small
fluctuations. Since $\bar\varphi$ satisfies the classical equations of
motion, the linear term vanishes and the quadratic Lagrangian is
\begin{equation}
    \mathcal{L}_{\text{quad}} = \frac{1}{2}(\partial_\mu\eta_i)^2
    - \frac{1}{2}\eta_i W^{ij}(\bar\varphi)\eta_j.
\end{equation}
If $W^{ij}$ is not diagonal, the fields $\eta_i$ are mixed and do not
correspond to independent physical particles. To decouple them, we
diagonalize $W$.

For an even quartic potential, the most general form is
\begin{equation}
    U(\varphi) = \frac{1}{2}m_{ij}\varphi_i\varphi_j
    + \frac{1}{4!}\lambda_{ijkl}\varphi_i\varphi_j\varphi_k\varphi_l,
\end{equation}
with $m_{ij} = m_{ji}$ and $\lambda_{ijkl} = \lambda_{(ijkl)}$.
Computing the second derivative,
\begin{equation}
    W^{ab}(\bar\varphi) = m_{ab}
    + \frac{1}{2}\lambda_{abkl}\bar\varphi_k\bar\varphi_l.
\end{equation}
Since $m_{ab}$ is symmetric and $\lambda_{abkl}$ is symmetric in $a,b$,
the matrix $W^{ab}$ is real and symmetric. By the spectral theorem, any
real symmetric matrix can be diagonalized by an orthogonal transformation:
there exists an orthogonal matrix $O$ such that
\begin{equation}
    O^T W O = \mathrm{diag}(w_1,\ldots,w_N).
\end{equation}
Defining the rotated fields $\chi_i = O_{ij}\eta_j$, and using the fact
that orthogonal transformations preserve the kinetic term, the quadratic
Lagrangian becomes
\begin{equation}
    \boxed{\mathcal{L}_{\text{quad}} = \sum_i\left[
    \frac{1}{2}(\partial_\mu\chi_i)^2
    - \frac{1}{2}w_i\,\chi_i^2\right].}
\end{equation}
Each mode $\chi_i$ is now an independent scalar with mass squared $w_i$,
propagating freely without mixing. These are the physical Higgs modes.

As a concrete example, for the gauge-invariant quartic potential
$U(\varphi) = \frac{1}{2}m^2\varphi_i\varphi_i
+ \frac{\lambda}{4!}(\varphi_i\varphi_i)^2$,
choosing the background along a single direction
$\bar\varphi_i = v\,\delta_{i1}$ gives
\begin{align}
    W^{11} &= m^2 + \frac{\lambda}{2}v^2
    \quad\text{(radial/Higgs mode)}, \\
    W^{aa} &= m^2 + \frac{\lambda}{6}v^2
    \quad\text{($N-1$ degenerate transverse modes, $a\neq 1$)}, \\
    W^{ab} &= 0 \quad\text{for $a\neq b$},
\end{align}
so the matrix is already diagonal in this basis.

\section{The Higgs Boson Contribution \texorpdfstring{$V_S^{(1)}$}{VS}}

For $V_S^{(1)}$, the only change relative to the scalar field theory
example is the presence of more than one scalar particle. The relevant
diagrams are those of the single scalar case, now with one internal line
associated with $\varphi_i$ and the other with $\varphi_j$.

% [DIAGRAM: one-loop scalar diagrams with external varphi lines,
%  analogous to Fig.~3 of Brandenberger but for the scalar loop]

The effective coupling constant, which includes the $\bar\varphi$ factors
for internal links and the Bose combinatorial factors, is
\begin{equation}
    W^{ij} = \frac{\partial^2 U(\varphi)}{\partial\varphi_i\,
    \partial\varphi_j}\bigg|_{\varphi=\bar\varphi}.
    \label{eq:Wij-coupling}
\end{equation}
In the single scalar field theory, $W = \frac{1}{2}\lambda\bar\varphi^2$.
If $U(\varphi)$ is an even quartic polynomial, $W(\bar\varphi)$ is a
quadratic form and can be diagonalized as shown above. In the diagonal
basis, the different Higgs fields decouple and $V_S^{(1)}$ is obtained
by adding the contributions of the new basis fields. Each individual term
is given by the scalar Coleman-Weinberg integral. Neglecting the
$\lambda^2$ corrections to terms already contained in $V_{\text{eff}}^{(0)}$,
the cutoff-independent part of the one-loop integral for a single mode
with $m^2 = \frac{\lambda}{2}\bar\varphi^2$ is
\begin{equation}
    V_{\text{1-loop}}\big|_{\text{c.i.}} = \frac{1}{64\pi^2}
    \left(\frac{\lambda\bar\varphi^2}{2}\right)^2
    \ln\frac{\lambda\bar\varphi^2}{2},
    \label{eq:CW-scalar}
\end{equation}
where the divergent terms proportional to $\Lambda^2$ and $\ln\Lambda^2$
have been absorbed into renormalization of the mass and coupling constant
respectively. Generalizing to the diagonal basis,
\begin{align}
    V_S^{(1)}(\bar\varphi) &= \frac{1}{64\pi^2}
    \sum_i(W^{ii})^2\ln W^{ii} \notag\\
    &= \frac{1}{64\pi^2}\mathrm{Tr}(W^2\ln W),
    \label{eq:VS1}
\end{align}
which is basis-independent since the trace is invariant under orthogonal
transformations.

\section{The Fermion Contribution \texorpdfstring{$V_f^{(1)}$}{Vf}}

The fermion contribution $V_f^{(1)}$ is given by the diagrams in Fig.~3
of Brandenberger, for the case $M=0$. Since the trace of an odd number
of $\gamma$ matrices vanishes, diagrams with an odd number of vertices
are zero. With all indices inserted explicitly, the Yukawa coupling term
from \eqref{eq:non-abelian-lagrangian} is
\begin{equation}
    \sum_{a,b,i}\bar\varphi_i\,\Gamma^i_{ab}\,\psi^a\psi^b.
    \label{eq:yukawa}
\end{equation}

% [DIAGRAM: closed fermion loop diagrams with external scalar insertions,
%  Fig.~3 of Brandenberger]

The diagram with two external scalar field insertions $\bar\varphi_i$
and $\bar\varphi_j$ consists of a closed fermion loop with two internal
massless propagators $\slashed{k}/k^2$ and two Yukawa vertices
contributing $\Gamma^i_{ab}$ and $\Gamma^j_{ba}$. Assembling all
factors, the overall sign $(-1)$ from the fermion loop and the symmetry
factor $\frac{1}{2}$ combine to give
\begin{equation}
    -\frac{1}{2}\int\frac{d^4k}{(2\pi)^4}
    \sum_{i,j,a,b}\mathrm{Tr}_S\,\bar\varphi_i\,\Gamma^i_{ab}\,
    \frac{\slashed{k}^2}{k^4}\,\Gamma^j_{ba}\,\bar\varphi_j.
    \label{eq:fermion-diagram}
\end{equation}
Using $(\slashed{k})(\slashed{k}) = k^2$, the propagator product
simplifies to $\slashed{k}^2/k^4 = 1/k^2$. The spinor trace then acts
only on the identity in spinor space, giving $\mathrm{Tr}_S(\mathbf{1})
= 4$ in four dimensions. Defining the matrix
$(\bar\varphi\,\Gamma)_{ab} \equiv \sum_i\bar\varphi_i\,\Gamma^i_{ab}$,
the sum over indices becomes
$\sum_{a,b}(\bar\varphi\,\Gamma)_{ab}(\bar\varphi\,\Gamma)_{ba}
= \mathrm{Tr}_f(\bar\varphi\,\Gamma)^2$,
and \eqref{eq:fermion-diagram} reduces to
\begin{equation}
    -\frac{1}{2}\int\frac{d^4k}{(2\pi)^4}\,\frac{4}{k^2}\,
    \mathrm{Tr}_f(\bar\varphi\,\Gamma)^2.
\end{equation}
The graph with $2n$ vertices yields analogously
\begin{equation}
    \frac{1}{2n}\int\frac{d^4k}{(2\pi)^4}\,\frac{1}{(k^2)^n}\,
    \mathrm{Tr}_f(\bar\varphi\,\Gamma)^{2n},
\end{equation}
so that summing over all diagrams,
\begin{equation}
    V_f^{(1)}(\bar\varphi) = -4i\,\mathrm{Tr}_f\sum_{n=1}^{\infty}
    \int\frac{d^4k}{(2\pi)^4}\,\frac{1}{2n}
    \left[\frac{(\bar\varphi\,\Gamma)^2}{k^2}\right]^n.
\end{equation}
This result is analogous to \eqref{eq:CW-scalar}, and its evaluation
follows the same procedure as in the scalar case. The cutoff-independent
part gives
\begin{equation}
    V_f^{(1)}(\bar\varphi) = -\frac{1}{64\pi^2}\,4\,
    \mathrm{Tr}_f\!\left[(\bar\varphi\,\Gamma)^4
    \ln(\bar\varphi\,\Gamma)^2\right].
    \label{eq:Vf1}
\end{equation}

\section{The Gauge Boson Contribution \texorpdfstring{$V_g^{(1)}$}{Vg}}

Finally we compute $V_g^{(1)}$, the gauge boson loop contribution. The
only vertex contributing to one-loop order is the $A^2\varphi^2$ vertex,
which stems from the following term in the Lagrangian density,
\begin{equation}
    \frac{1}{2}\mathrm{Tr}(igA_\mu\varphi)^\dagger(igA^\mu\varphi)
    = \frac{1}{2}g_{\alpha\beta}\mathrm{Tr}
    \!\left[(\tau_\alpha A^\alpha_\mu\varphi)^\dagger
    (\tau_\beta A^{\beta\mu}\varphi)\right]
    = \frac{1}{2}A^\alpha_\mu M^2_{\alpha\beta}A^{\beta\mu},
\end{equation}
with the field-dependent gauge boson mass matrix
\begin{equation}
    M^2_{\alpha\beta} = g_{\alpha\beta}\,
    \mathrm{Tr}\!\left[(\tau_\alpha\varphi)^\dagger\tau_\beta\varphi\right].
    \label{eq:gauge-mass-matrix}
\end{equation}

% [DIAGRAM: one-loop gauge boson diagrams with external scalar
%  insertions, Fig.~4 of Brandenberger]

A few comments on \eqref{eq:gauge-mass-matrix} are in order. First, if
$G$ is not a simple group, there will in general be a different coupling
constant associated with each field $A^a_\mu$. Second, the trace is in
the Higgs representation space. Third, $M^2_{\alpha\beta}|_{\varphi=
\bar\varphi}$ is the gauge boson mass matrix in the broken phase with
vacuum expectation value $\bar\varphi$, and is a constant matrix in the
sense of \eqref{eq:Wij-def}. Since in Landau gauge the gauge fields have
three degrees of freedom (they have a Lie algebra trace), the calculation
of $V_g^{(1)}$ is analogous to that of $V_S^{(1)}$, and
\begin{equation}
    V_g^{(1)}(\bar\varphi) = \frac{3}{64\pi^2}
    \mathrm{Tr}(M^4\ln M^2)\big|_{\varphi=\bar\varphi}.
    \label{eq:Vg1}
\end{equation}

\section{The Full One-Loop Effective Potential}

Combining the three contributions \eqref{eq:VS1}, \eqref{eq:Vf1}, and
\eqref{eq:Vg1}, the full one-loop effective potential is
\begin{equation}
    V_{\text{eff}}(\bar\varphi) = U(\bar\varphi)
    + \frac{1}{64\pi^2}\mathrm{Tr}(W^2\ln W)
    - \frac{4}{64\pi^2}\mathrm{Tr}_f
    \!\left[(\bar\varphi\,\Gamma)^4\ln(\bar\varphi\,\Gamma)^2\right]
    + \frac{3}{64\pi^2}\mathrm{Tr}(M^4\ln M^2)\big|_{\varphi=\bar\varphi},
\end{equation}
where all traces are evaluated at $\varphi = \bar\varphi$. This is the
master formula for the one-loop effective potential in a non-Abelian
gauge theory. The scalar, fermion, and gauge boson loops each contribute
a term of the Coleman-Weinberg form $m^4\ln m^2$, with the mass matrix
appropriate to each field type. In the next chapter we apply this formula
to GUT models, where the gauge group is larger and the symmetry breaking
pattern is richer.
